% ================================================================
%  DROP-IN REPLACEMENT — paste over the two existing reeds blocks
%  1. \subsection{Bushes/Shrubs experiments - Reeds}   (Methodology)
%  2. \subsection{Reeds/bushes scan}                   (Testing and Results)
% ================================================================


% ----------------------------------------------------------------
%  METHODOLOGY
% ----------------------------------------------------------------
\subsection{Bushes/Shrubs experiments - Reeds}

Two radar scans of the reed bed were acquired to test detection performance
under different levels of foliage occlusion.
The first scan was conducted on 11~March~2026, with the corner reflectors
elevated on stands so that they were only partially occluded by the upper
reed stems.
The second scan was conducted on 8~April~2026, with the corner reflectors
placed directly at ground level behind the reeds, resulting in heavy occlusion
with the reflectors barely visible to the naked eye.
The same radar scan parameters were used on both days to allow a direct
comparison of how occlusion severity affects detection performance.

\subsubsection{Scan Procedure for Reeds}

The scan parameters used on both days are given in Table~\ref{tab:scan_parameters}.
Both scans were conducted under sunny, slightly windy conditions at approximately
27$^\circ$C with no cloud cover.
On both days, two corner reflectors were placed side by side behind the reeds,
with the radar positioned at the same location for consistency.

In the March scan, the reflectors were elevated on stands above the densest part
of the reed stems, reducing the effective path length through the vegetation.
In the April scan, the reflectors were placed on gravel at ground level,
maximising the path length through the full reed height.
The trihedral corner reflectors were selected as the detection target because
they have a high radar cross-section (RCS) and strong geometric reflectivity,
producing a distinct high-intensity return in both the radar and LiDAR datasets.
This provides a controlled reference target for assessing concealment effects and
directly comparing reflector detectability across sensing modalities.

\subsubsection{Behind-Foliage Region Definition}

The detection analysis was spatially constrained to the region between the front
face of the reeds and the corner reflector.
This definition was chosen because it includes the vegetation returns as part of
the clutter, which gives a more realistic false positive rate than using only the
narrow gap immediately behind the foliage.
Using the back face of the reeds as the boundary was considered but rejected, as
the resulting region contained only 8--10 points, making any false positive rate
statistically meaningless.
The reed front face was identified at a range of 4.1~m in March and 3.0~m in
April from the radar.
All returns in this window were included in the clutter floor and false positive
calculations.

\subsubsection{Detection Threshold Approaches}

Two complementary threshold approaches were used to analyse detection performance,
each serving a different analytical purpose.

The first approach used a fixed threshold derived from the control scan clutter
floor ($-40.1$~dB).
Applying the same threshold to both foliage scans enables direct comparison
between March and April on a common reference, and allows vegetation loss to be
computed relative to free-space conditions.

The second approach estimated the clutter floor independently for each scan from
the returns within the behind-foliage region.
This is more representative of how a threshold would be set in a real deployment,
where no prior free-space reference is available.
The trade-off is that it makes direct comparison between scans harder, since each
scan operates against a different baseline.

Both approaches detected the CR successfully in all scans.
The fixed threshold approach is used for the vegetation loss calculations, and the
scan-specific approach is used for the ROC curve comparison.

\subsubsection{Vegetation Loss Calculation}

Vegetation loss was estimated by comparing the measured CR peak amplitude in each
foliage scan against the expected free-space amplitude at the same range, derived
from the control scan propagation model.
The vegetation loss is then:
\[
  L_{\text{veg}} = A_{\text{free-space}} - A_{\text{measured}}
\]
where $A_{\text{free-space}}$ is the predicted amplitude in the absence of
vegetation and $A_{\text{measured}}$ is the actual peak amplitude observed in the
foliage scan.


% ----------------------------------------------------------------
%  TESTING AND RESULTS
% ----------------------------------------------------------------
\subsection{Reeds/Bushes Scan}

\subsubsection{Radar Scans}

Two scans of the reeds were taken — one on the 11th March 2026 at 11:48 and one
on the 8th April 2026 at 13:19.
Two individual corner reflectors were placed behind the reeds on each day, side by
side.
On both days, the same radar scan parameters were used for consistency, as shown
in Table~\ref{tab:scan_parameters}.

\begin{table}[H]
    \centering
    \begin{tabular}{lccc}
        \hline
        \textbf{Parameter} & \textbf{Min} & \textbf{Max} & \textbf{Value} \\
        \hline
        spacing    & 0.1    & 10.0  & 0.5    \\
        min\_pitch & -100.0 & -1.0  & -100.0 \\
        max\_pitch & -100.0 & -1.0  & -60.0  \\
        min\_roll  & 0.0    & 360.0 & 0.0    \\
        max\_roll  & 0.0    & 360.0 & 40.0   \\
        \hline
    \end{tabular}
    \caption{Radar scan parameters}
    \label{tab:scan_parameters}
\end{table}

\paragraph{Scan 1: 11/03/2026}
Sunny conditions, no clouds, 27$^\circ$C, slightly windy.
On the 11th March, the corner reflectors were placed much higher up and were only
partially obscured by the reeds.
On this scan both corner reflectors are seen clearly on the radar scans.
The setup is shown in Figure~\ref{fig:cr_reeds}.

\begin{figure}[H]
    \centering
    \begin{subfigure}{0.48\linewidth}
        \includegraphics[width=\linewidth]{images/reeds_scan/cr_through_reeds_higher.jpeg}
        \caption{Enter Caption}
        \label{fig:cr_higher}
    \end{subfigure}
    \hfill
    \begin{subfigure}{0.48\linewidth}
        \includegraphics[width=\linewidth]{images/reeds_scan/zoomed_in_reeds_cr.png}
        \caption{Enter Caption}
        \label{fig:seethrough}
    \end{subfigure}
    \caption{Reeds scan setup --- 11th March 2026}
    \label{fig:cr_reeds}
\end{figure}

\paragraph{Scan 2: 8/04/2026}
Sunny conditions, no clouds, 27$^\circ$C, slightly windy.
On the 8th of April, the corner reflectors were placed much lower to the ground,
meaning they were obscured by the reeds and hardly visible to the naked eye.
As can be seen in Figure~\ref{fig:reeds_april_setup}, only one corner reflector
is visible to the human eye — the smaller one to the left cannot be seen at all.

\begin{figure}[H]
    \centering
    \begin{subfigure}[b]{0.48\textwidth}
        \centering
        \includegraphics[width=\linewidth]{images/reeds_scan/reeds_camera_image_11th_march.jpeg}
        \caption{Radar scanning reeds, 8th April}
        \label{fig:reeds_march}
    \end{subfigure}
    \hfill
    \begin{subfigure}[b]{0.48\textwidth}
        \centering
        \includegraphics[width=\linewidth]{images/reeds_scan/reeds_camera_image_11th_march_cr.jpeg}
        \caption{Corner reflector placed low to the ground}
        \label{fig:reeds_march_cr}
    \end{subfigure}
    \caption{Reeds scan setup --- 8th April 2026}
    \label{fig:reeds_april_setup}
\end{figure}

\subsubsection{Point Cloud Visualisation}

Figure~\ref{fig:reeds_radar_comparison} shows the 3D radar point clouds for the
March and April scans coloured by amplitude.

\begin{figure}[H]
    \centering
    \begin{subfigure}[b]{0.48\textwidth}
        \centering
        \includegraphics[width=\linewidth]{images/reeds_scan/march_reeds_radar.png}
        \caption{March reeds scan coloured by amplitude}
        \label{fig:march_reeds_radar}
    \end{subfigure}
    \hfill
    \begin{subfigure}[b]{0.48\textwidth}
        \centering
        \includegraphics[width=\linewidth]{images/reeds_scan/april_reeds_radar.png}
        \caption{April reeds scan coloured by amplitude}
        \label{fig:april_reeds_radar}
    \end{subfigure}
    \caption{3D radar point cloud of reeds scans with corner reflector locations marked}
    \label{fig:reeds_radar_comparison}
\end{figure}

The dominant factor driving the difference in return amplitude between the two
scans is the height of the corner reflectors relative to the reed canopy.
There is approximately a 64~dB difference between corner reflector 2 from March
to April, even though the reeds themselves have not changed between scans.

A secondary factor is the surface material beneath the corner reflectors.
In April, the corner reflectors were placed on gravel, which can create multipath
where the radar signal bounces off the ground and interferes constructively with
the direct CR return — this effect is much stronger at ground level.
In March, the corner reflectors were placed on grass, which absorbs more energy
and produces less multipath.

\subsubsection{Detection Analysis}

\paragraph{Scan-specific thresholds.}

Unlike the fixed control threshold approach, here the clutter floor is estimated
directly from each scan's own behind-foliage returns, which is more representative
of a real deployment where no prior free-space reference is available.
The March clutter floor was $-31.8$~dB and the April clutter floor was $-33.0$~dB.
At these thresholds, CR1 was detected in both scans while April CR2 remained
undetectable, falling below its own scan's clutter floor at $-11.0$~dB.
The results are consistent with the fixed threshold analysis above.
Figures~\ref{fig:reeds_amplitude_vs_range}--\ref{fig:scanspec_table} show the
amplitude vs range, threshold sweep, ROC curves, and detection summary for both
scans with CR1 and CR2 reported separately.

\begin{figure}[H]
    \centering
    \includegraphics[width=\linewidth]{figures/scanspecific/fig1_scanspec_amplitude_vs_range.png}
    \caption{Amplitude vs range for the March (left) and April (right) reeds
    scans using scan-specific clutter floors. CR1 and CR2 are labelled separately.}
    \label{fig:reeds_amplitude_vs_range}
\end{figure}

\begin{figure}[H]
    \centering
    \includegraphics[width=\linewidth]{figures/scanspecific/fig3_scanspec_threshold_sweep.png}
    \caption{CR points detected (left) and false positives (right) as a function
    of threshold for both scans, with CR1 and CR2 shown separately.}
    \label{fig:scanspec_threshold_sweep}
\end{figure}

\begin{figure}[H]
    \centering
    \includegraphics[width=0.75\linewidth]{figures/scanspecific/fig2_scanspec_roc_curves.png}
    \caption{ROC curve comparison using scan-specific thresholds. Operating
    points are marked at each scan's own clutter floor.}
    \label{fig:scanspec_roc}
\end{figure}

Table~\ref{fig:scanspec_table} summarises the scan-specific detection metrics
for CR1 and CR2 across both scans.

\begin{figure}[H]
    \centering
    \includegraphics[width=\linewidth]{figures/scanspecific/fig4_scanspec_summary_table.png}
    \caption{Scan-specific detection summary for CR1 and CR2 across the March
    and April reeds scans.}
    \label{fig:scanspec_table}
\end{figure}

\paragraph{Fixed control threshold.}

At the fixed control threshold of $-40.1$~dB, the March scan achieved a true
positive rate (TPR) of 0.93 and a false positive rate (FPR) of 0.67.
The April scan achieved a TPR of 0.57 with the same FPR of 0.67.
The identical FPR across both scans indicates that the behind-foliage clutter
distribution was similar in both sessions, meaning the vegetation structure did
not change significantly between March and April.
The lower TPR for April is consistent with heavier occlusion reducing the number
of detectable CR returns above the threshold.

Figure~\ref{fig:reeds_snr_comparison} shows the SNR for each individual corner
reflector across both scans, with April CR2 marked as undetected.

\begin{figure}[H]
    \centering
    \includegraphics[width=0.75\linewidth]{figures/fig4_reeds_snr_comparison.png}
    \caption{SNR above clutter floor for each corner reflector in the March and
    April reeds scans. April CR2 is marked as undetected as it fell below the
    clutter floor.}
    \label{fig:reeds_snr_comparison}
\end{figure}

Figure~\ref{fig:reeds_detection_summary} summarises the detection metrics for
each corner reflector in both scans.

\begin{figure}[H]
    \centering
    \includegraphics[width=\linewidth]{figures/fig5_reeds_detection_summary_table.png}
    \caption{Detection summary for CR1 and CR2 across the March and April reeds
    scans. April CR2 is highlighted as undetected.}
    \label{fig:reeds_detection_summary}
\end{figure}

\subsubsection{Vegetation Attenuation Estimate}

Vegetation loss was estimated by comparing the measured CR peak amplitude in each
foliage scan against the expected free-space amplitude at the same range, derived
from the control scan.

For the March scan, the expected free-space amplitude at 10.92~m was 76.3~dB.
The measured CR peak was 18.0~dB, giving a vegetation loss of:
\[
  L_{\text{March}} = 76.3 - 18.0 = 58.3~\text{dB}
\]

For the April scan, the expected free-space amplitude at 13.88~m was 72.2~dB.
The measured CR1 peak was 53.0~dB, giving a vegetation loss of:
\[
  L_{\text{April}} = 72.2 - 53.0 = 19.2~\text{dB}
\]

Figure~\ref{fig:reeds_attenuation_vs_range} shows the peak amplitude of each
corner reflector plotted against its range.
Notably, April CR1 is further from the sensor (13.88~m) than the March
reflectors (10.92~m) yet produces a significantly stronger return, demonstrating
that target height and alignment relative to the vegetation dominates over range
as the primary factor in detection performance.

\begin{figure}[H]
    \centering
    \includegraphics[width=0.75\linewidth]{images/reeds_scan/reeds_attenuation_vs_range.png}
    \caption{CR peak amplitude vs range for March and April reeds scans.
    The dashed line shows the detection threshold ($-29.2$~dB).}
    \label{fig:reeds_attenuation_vs_range}
\end{figure}

The April vegetation loss (19.2~dB) is considerably lower than March (58.3~dB),
despite the April reflectors being more heavily occluded.
This is counterintuitive and is largely attributed to multipath from the gravel
surface beneath the April reflectors, as discussed in the point cloud section
above.
The April estimate of 19.2~dB should therefore be treated as a lower bound on
the true path loss through the reeds at ground level.

\subsubsection{Amplitude Distributions}

Figure~\ref{fig:reeds_amplitude_hist} shows the amplitude distributions for both
scans.

\begin{figure}[H]
    \centering
    \includegraphics[width=\linewidth]{images/reeds_scan/reeds_amplitude_histograms.png}
    \caption{Amplitude distributions for the March (left) and April (right) reeds
    scans. Clutter floor, detection threshold, and CR peak values are marked.}
    \label{fig:reeds_amplitude_hist}
\end{figure}

The March distribution is broad and relatively flat, centred around $-25$~dB,
with the CR peaks at approximately 18~dB appearing as a clear spike at the far
right — well separated from the vegetation clutter.
The April distribution is more compact, centred around $-28$~dB.
April CR2 ($-11.0$~dB) sits within the upper tail of the clutter distribution,
explaining why it is undetectable above the threshold.
April CR1 (53.0~dB) is an extreme outlier far beyond the clutter body, consistent
with the multipath-boosted return discussed above.
The broader March distribution reflects greater variability in returns from
partially-occluded vegetation, while the more compact April distribution suggests
a more uniform clutter background at ground level.

\subsubsection{LiDAR Scans}

\begin{figure}[H]
    \centering
    \includegraphics[width=1\linewidth]{images/reeds_lidar_scan.png}
    \caption{LiDAR scan of reeds in Cadigal Greens with corner reflector behind}
    \label{fig:reeds_lidar}
\end{figure}

\begin{figure}[h]
    \centering
    \begin{subfigure}{0.48\linewidth}
        \includegraphics[width=\linewidth]{images/reeds_scan/reeds_lidar_by_height.png}
        \caption{Enter Caption}
        \label{fig:height}
    \end{subfigure}
    \hfill
    \begin{subfigure}{0.48\linewidth}
        \includegraphics[width=\linewidth]{images/reeds_scan/reeds_lidar_by_reflectivity.png}
        \caption{Enter Caption}
        \label{fig:reflectivity}
    \end{subfigure}
    \caption{Overall Caption}
    \label{fig:both}
\end{figure}

Figure~\ref{fig:number_vs_distance_lidar} shows that the LiDAR produces a high
density of surface returns from the front face of the reed bed, however the point
count drops sharply beyond this surface, indicating that the laser pulses cannot
penetrate the vegetation.
As a result, any corner reflectors placed within or behind the reeds are
completely obscured and produce no detectable return.

\begin{figure}[H]
    \centering
    \includegraphics[width=0.75\linewidth]{images/reeds_scan/number_vs_distance_lidar.png}
    \caption{LiDAR point return count as a function of distance from the sensor.}
    \label{fig:number_vs_distance_lidar}
\end{figure}

This is in direct contrast to the radar, which successfully detected the CR
through the reeds in all tested configurations.
The result demonstrates the fundamental advantage of millimetre-wave radar over
optical LiDAR for concealed-target detection in dense vegetation.
